\documentclass[14pt,a4paper]{extarticle} 
\usepackage[T2A]{fontenc}
\usepackage[utf8]{inputenc}
\usepackage[top=0.78in, bottom=0.78in, right=0.59in, includefoot]{geometry}
\usepackage[english, russian]{babel} % Подключение русского языка
\usepackage{indentfirst} % Абзацы
\usepackage{amsmath, amsfonts, amssymb, amsthm, mathtools}
\usepackage[T2A]{fontenc}
\usepackage[utf8]{inputenc}

\usepackage{titlesec}
% Центрирование \section
\titleformat{\section} % Команда для настройки \section
{\centering\normalfont\Large\bfseries} % Стиль: центрирование, крупный шрифт, жирный
{\thesection} % Номер раздела
{1em} % Расстояние между номером и заголовком
{} % Дополнительный код после заголовка

\usepackage[hidelinks]{hyperref} % подключение ссылок
\usepackage{nameref} 
\usepackage{csquotes} 

\usepackage{hyphsubst} % Этот пакет позволяет настраивать правила переноса слов
\usepackage{minted} % Используется для форматирования и подсветки синтаксиса кода

\usepackage{titleps} % создание линий для колонтитулов
\newpagestyle{main}{
	% \setfootrule{0.4pt}
	\setfoot{}{\thepage}{}
}

\pagestyle{main}

%\linespread{1} % Отступы между строками
%\setlength{\parindent}{20pt}
%\setlength{\parskip}{10pt}

% Создание теорем, определений и прочего. Можно поменять стиль, сделать разное оформление к каждому разделу, но придется дописывать преамбулу
\usepackage{amsthm}
\theoremstyle{plain}
\newtheorem{theorem}{Теорема}
\newtheorem*{definition}{Определение}
\newtheorem*{corollary}{Следствие}
\newtheorem*{note}{Замечание}


\usepackage{tocloft} % Позволяет настраивать внешний вид оглавления (table of contents), списка фигур (list of figures) и списка таблиц (list of tables)
\renewcommand{\cftsecleader}{\cftdotfill{\cftdotsep}} % Заполнение пробелов между заголовками секций и номерами страниц в оглавлении
%\renewcommand{\thesection}{\arabic{section}.} % Переопределяет формат отображения номера секции. В данном случае, номера секций будут отображаться в арабском формате (1, 2, 3 и т.д.) и будут следовать с точкой (используйте, если у вас нет поздразделов, иначе у них будет писать 1..1 и так далее)

%Создание изображений
\usepackage{graphicx} % Вставка изображений
%\usepackage{float} % "Плавающие" картинки
\usepackage{wrapfig} % Обтекание фигур (таблиц, картинок и прочего)
\usepackage{placeins} % Этот пакет управляет размещением плавающих объектов (таких как рисунки и таблицы) в документе
\usepackage[dvipsnames]{xcolor} % Предоставляет расширенные возможности работы с цветами
\usepackage{caption} % Используется для настройки внешнего вида подписей к рисункам и таблицам
\captionsetup[figure]{name=Рисунок}
\DeclareCaptionLabelSeparator{dash}{ --- } % Эта команда определяет разделитель между названием (например, "Рисунок") и текстом подписи
\captionsetup{labelsep=dash} % Эта команда применяет ранее определенный разделитель (в данном случае dash) ко всем подписям в документе

% Для таблиц
\usepackage[raggedrightboxes]{ragged2e}
\usepackage{makecell} 
\usepackage{tabularx}
\usepackage{enumitem} % Пакет для гибкой настройки списка
\usepackage{multicol}
\usepackage{multirow} 
\usepackage{booktabs}
\usepackage{float}
\usepackage{array}  
\usepackage{longtable}
\usepackage{pdflscape}
% Создание листинга программы
\usepackage{listings}
\usepackage{xcolor}


\lstset{
	language=C++,
	basicstyle=\ttfamily\footnotesize,
	keywordstyle=\color{blue}\bfseries,
	commentstyle=\color{green!60!black},
	stringstyle=\color{red},
	numbers=left,
	frame=none,
}

\newcommand{\mytitlepage}{
	\thispagestyle{empty}
	
	\centerline{\textbf{Министерство науки и высшего образования РФ}}
	\centerline{\textbf{Федеральное государственное бюджетное образовательное}}
	\centerline{\textbf{учреждение высшего образования}}
	\centerline{\textbf{<<Уфимский университет науки и технологий>>}}
	\vspace{5mm}
	\centerline{\textbf{Кафедра:} Высокопроизводительных вычислений и дифференциальных}
	\centerline{уравнений}
	
	
	\vspace{40mm}
	
	\centerline{\textbf{\makecell{ИССЛЕДОВАНИЕ ЭВОЛЮЦИИ НЕЛИНЕЙНОЙ ДИССИПАТИВНОЙ
				\\ДИНАМИЧЕСКОЙ СИСТЕМЫ }}} % Написать все название капсом
	\vspace{5mm}
	\centerline{{\textbf{ОТЧЕТ}}}
	\vspace{2mm}
	\centerline{к лабораторной работе №2 по дисциплине}
	\vspace{2mm}
	\centerline{<<Математическое моделирование>>} % Написать название дисциплины
	%\vspace{4mm}
	%\centerline{\textbf{3952.236108.000 ПЗ}} % Посередине шифр ABCDEF
	% A - тип работы, где 2 - курсовой проект, 3 - курсовая работа (оставляем 2)
	% BС - код дисциплины, (состоит из кода учебной программы 3 и номера семестра)
	% D - обозначение группы: 1 у ПМ, 2 у Мкн
	% EF - номер в списке группы по журналу
	\vspace{30mm}
	% Указать группу, прописать фамилии и инициалы в таблице
	\begin{center}
		\bgroup
		\def\arraystretch{1.4}
		\setlength{\tabcolsep}{12pt}
		\begin{tabular}{|c|c|c|c|c|}
			\hline
			\makecell{Группа\\МКН-418Б} & Фамилия И.\ О. & Подпись & Дата & Оценка \\
			\hline
			Студент & Халитова А.А. &  &  & \\
			\hline
			Принял & Феоктистов Б.А. &  &  & \\
			\hline
		\end{tabular}
		\egroup
	\end{center}
	\vspace{35mm}
	\centerline{Уфа 2026} % Поменять год 
	\clearpage
}

\begin{document}
	\mytitlepage
	\tableofcontents
	\newpage
	
	\section{Введение}
	
	\textbf{Цель работы:} получить навык численного исследования динамики нелинейной
	диссипативной динамической системы, обладающей странным аттрактором.
	
	\newpage
	
	\section{Теоретическая часть}
	
	Рассматривается нелинейная двухпараметрическая автономная динамическая система 
	
	\begin{equation}
		\begin{cases}
			\dfrac{dx}{dt} = f(x,y,z,a,b), \\[12pt]
			\dfrac{dy}{dt} = g(x,y,z,a,b), \\[12pt]
			\dfrac{dz}{dt} = h(x,y,z,a,b).
		\end{cases} \nonumber
	\end{equation}
	
	Функции f, g, h выбираются в соответствии с номером варианта. Для заданной системы выполнить следующие задания:
	
	\begin{enumerate}
		\item Определить области изменения параметров a и b, в которых данная динамическая система является диссипативной.
		\item Определить стационарные точки диссипативной системы.
		\item Исследовать стационарные точки на асимптотическую устойчивость по
		первому приближению.
		\item Определить значения параметров a и b, при которых в системе появляется странный аттрактор.
		\item Написать вычислительную программу на языке программирования C++,
		реализующую процедуру численного интегрирования исходной диссипативной системы по методу Рунге-Кутта 4-го порядка точности.
		\item С использованием вычислительной программы провести серию вычислительных экспериментов, демонстрирующих различные виды динамики системы. Построить траектории системы в окрестности стационарных точек. Определить численно значения параметров a и b, при которых в системе существует странный аттрактор и при которых система переходит в режим автоколебаний. 
	\end{enumerate}
	
	\newpage
	
	\section{Практическая часть}
	
	Система уравнений для 7 варианта:
	
	\begin{equation}
		\begin{cases}
			\dfrac{dx}{dt} = a(y-x) = f_x, \\[12pt]
			\dfrac{dy}{dt} = bx + y - xz = g_y, \\[12pt]
			\dfrac{dz}{dt} = -\dfrac{8}{3}z + xy = h_z.
		\end{cases}
			\label{sys:main_equations}
	\end{equation}
	
	\subsection{Область определения диссипативной системы}
	
	Чтобы система была диссипативной, необходимо выполнение условия:
	
	\begin{equation}
		\vec{q} = (f_x, g_y, h_z)^T, \quad \operatorname{div} \vec{q} < 0.
		\label{sys:dissipative_conditions}
	\end{equation}
	
	После подстановки уравнений \eqref{sys:main_equations} в \eqref{sys:dissipative_conditions}, получим:
	
	\begin{equation}
		\operatorname{div} \vec{q} = -a - \frac{5}{3}, \quad \Rightarrow \quad a > -\frac{5}{3}.\nonumber
	\end{equation}
	
	\subsection{Стационарные точки системы}
	
	\textbf{Стационарные точки} системы находятся из решения системы уравнений:
	
	\begin{equation}
		\begin{cases}
			a(y-x) = 0, \\[12pt]
			bx + y - xz = 0, \\[12pt]
			-\dfrac{8}{3}z + xy = 0.
		\end{cases} \nonumber
	\end{equation}
	
	Получаем следующие стационарные точки:
	
	\begin{equation}
		O_1 = (0, 0, 0),\nonumber
	\end{equation}
	\begin{equation}
		O_2 = \left( \sqrt{\frac{8}{3}(b+1)}; \sqrt{\frac{8}{3}(b+1)}; b+1 \right),\nonumber
	\end{equation}
	\begin{equation}
		O_3 = \left( -\sqrt{\frac{8}{3}(b+1)}; -\sqrt{\frac{8}{3}(b+1)}; b+1 \right),\nonumber
	\end{equation}
	
	с дополнительным ограничением $b > -1$.
	
	\subsection{Исследование на асимптотическую устойчивость по первому приближению}
	
	Условия устойчивости системы:
	
	\begin{equation}
		\operatorname{Re}(\lambda_{i})<0.\nonumber
	\end{equation}
	
	Сделаем замену переменных:
	\begin{equation}
		\centering
		\begin{cases}
			x = x_k + \xi, \\[12pt]
			y = y_k + \eta, \\[12pt]
			z = z_k + \zeta,
		\end{cases}
		\quad \Rightarrow \quad
		\begin{cases}
			\dfrac{d\xi}{dt}=a((y_{k}+\eta)-(x_{k}+\xi)),\\[12pt]
			\dfrac{d\eta}{dt}=b(x_{k}+\xi)+(y_{k}+\eta)-(x_{k}+\xi)(z_{k}+\zeta),\\[12pt]
			\dfrac{d\zeta}{dt}=-\dfrac{8}{3}(z_{k}+\zeta)+(x_{k}+\xi)(y_{k}+\eta),
		\end{cases}\nonumber
		\label{eq:linearization}
	\end{equation}
	где $x_{k},y_{k},z_{k}$ - координаты стационарных точек, $\xi,\eta,\zeta$ - малые отклонения.
	
	Разложим в ряд Тейлора в окрестностях точек:
	
	\begin{equation}
		\begin{cases}
			\dfrac{d\xi}{dt} = -a\xi + a\eta, \\[12pt]
			\dfrac{d\eta}{dt} = (b - z_k)\xi + \eta - x_k\zeta, \\[12pt]
			\dfrac{d\zeta}{dt} = y_k\xi + x_k\eta - \dfrac{8}{3}\zeta.
		\end{cases}
		\nonumber
	\end{equation}
	
	Получим матрицу правых частей:
	
	\begin{equation}
		M = \begin{pmatrix}
			-a & a & 0 \\[6pt]
			b - z_k & 1 & -x_k \\[6pt]
			y_k & x_k & -\dfrac{8}{3}
			\end{pmatrix}.
		\label{sys:right_part_matrix}
	\end{equation}
	
	Рассмотрим точку $O_1$, для которой матрица ~\ref{sys:right_part_matrix} будет иметь вид:
	
	\begin{equation}
		M = \begin{pmatrix}
			-a & a & 0 \\[6pt]
			b  & 1 & 0 \\[6pt]
			0 & 0 & -\dfrac{8}{3}
		\end{pmatrix}.\nonumber
	\end{equation}
	
	Ее характеристическое уравнение будет иметь вид:
	
	\begin{equation}
		\lambda^3 + \lambda^2\left(\frac{5}{3} + a\right) + \lambda\left(\frac{5}{3}a - \frac{8}{3} + ab\right) - \frac{8}{3}a(b+1) = 0.
		\label{eq:char_eq}\nonumber
	\end{equation}
	
	По критерию Гурвица, необходимым и достаточным условием отрицательности действительных частей всех корней характеристического уравнения является положительность всех главных диагональных миноров матрицы Гурвица.
	
	Для характеристического уравнения третьего порядка общего вида:
	
	\begin{equation}
		\lambda^3 + p_1\lambda^2 + p_2\lambda + p_3 = 0,\nonumber
	\end{equation}
	матрица Гурвица имеет вид:
	
	\begin{equation}
		H = \begin{pmatrix}
			p_1 & p_0 & 0 \\
			p_3 & p_2 & p_1 \\
			p_5 & p_4 & p_3
		\end{pmatrix}.\nonumber
	\end{equation}
	
	Получим условия устойчивости:
	
	\begin{equation}
		\begin{cases}
			\Delta_1 = \dfrac{5}{3} + a > 0 \quad \Rightarrow \quad a > -\dfrac{5}{3}, \\[12pt]
			\Delta_2 = \dfrac{8}{9}\left(3a^2 + 6ab + 5a - 10\right) > 0, \\[12pt]
			\Delta_3 = \dfrac{8}{27}a(b+1)\left(9a^2 + 9ab - 12a - 20b - 20\right) > 0.
		\end{cases}\nonumber
	\end{equation}
	
	Тогда чтобы точка $O_1 = (0,0,0)$ была устойчива, необходимо выполнение условий:
		
	\begin{equation}
		\begin{cases}
			a > 1, \\[12pt]
			b > -1.
		\end{cases}
		\label{conditions_1_point}
	\end{equation}
	
	Область ~\eqref{conditions_1_point} показана на рисунке ~\ref{fig:rightpart}:
	
	\begin{figure}[H]
		\centering
		\includegraphics[width=0.6\linewidth]{../../../../../../Pictures/Screenshots/right_part}
		\caption{Область устойчивости точки $O_1$}
		\label{fig:rightpart}
	\end{figure}
	
	Рассмотрим точку $O_2 = \left( \sqrt{\frac{8}{3}(b+1)}; \sqrt{\frac{8}{3}(b+1)}; b+1 \right)$. Характеристическое уравнение для этой точки будет иметь вид:
	
	\begin{equation}
		\lambda^3+\lambda^2\left(\frac{5}{3}+a\right)+\lambda\frac{8}{3}\left(a+b\right)+\frac{16}{3}a(b+1)=0.\nonumber
	\end{equation}
	
	Получим условия устойчивости:
	
	\begin{equation}
		\begin{cases}
			\Delta_1 = \dfrac{5}{3} + a > 0 \quad \Rightarrow \quad a > -\dfrac{5}{3}, \\[12pt]
			\Delta_2 = \frac{8}{9}\left(3a^2 - 3ab - a + 5b\right) > 0,\\[12pt]
			\Delta_3 = \frac{128}{27}a(b+1)\left(3a^2 - 3ab - a + 5b\right) > 0.
		\end{cases}\nonumber
	\end{equation}
	
	Тогда чтобы точка $O_2$ была устойчива, необходимо выполнение условий (рисунок ~\ref{fig:secpndpart}):
	
	\begin{equation}
		\begin{cases}
			a \in [0,1], & b \in \left(\dfrac{a-3a^2}{5-3a}, +\infty\right), \\[12pt]
			a \in \left(1,\dfrac{5}{3}\right), & b \in (-1, +\infty), \\[12pt]
			a \in \left(\dfrac{5}{3}, + \infty\right), & b \in \left(-1, \dfrac{a-3a^2}{5-3a}\right).
		\end{cases}
		\nonumber
	\end{equation}
	
	\begin{figure}[H]
		\centering
		\includegraphics[width=0.6\linewidth]{../../../../../../Pictures/Screenshots/secpnd_part}
		\caption{Область устойчивости точки $O_2$}
		\label{fig:secpndpart}
	\end{figure}

	\subsection{Определение параметров возникновения странного аттрактора}
	
	В ходе вычислительных экспериментов были исследованы различные режимы динамики системы в зависимости от параметров $a$ и $b$. Начальные условия задавались вблизи соответствующих стационарных точек с малым возмущением $\delta = 0.1$.
	
	\subsubsection{Исследование точки $O_1$}
	
	\begin{figure}[H]
		\centering
		\includegraphics[width=0.7\linewidth]{../../../build/lr2_screen/plot_a2.0_b-2.0}
		\caption{Устойчивая точка $O_1$. Точки $O_2$ и $O_3$ не существуют}
		\label{fig:O1_stable}
	\end{figure}
	
	\begin{figure}[H]
		\centering
		\includegraphics[width=0.7\linewidth]{../../../build/lr2_screen/plot_a0.5_b5.0}
		\caption{Неустойчивая точка $O_1$. Точки $O_2$ и $O_3$ устойчивы}
		\label{fig:O1_unstable_O2_stable}
	\end{figure}
	
	На рисунке~\ref{fig:O1_unstable_O2_stable} видно, что траектория, начиная вблизи точки $O_1$,  сходится к устойчивой точке $O_2$.
	
	\begin{figure}[H]
		\centering
		\includegraphics[width=0.6\linewidth]{../../../build/lr2_screen/plot_a1.2_b2.0}
		\caption{Все точки устойчивы}
		\label{fig:all_stable}
	\end{figure}
	
	На рисунке~\ref{fig:boundary} траектория демонстрирует длительный переходной процесс, что характерно для бифуркационных значений параметров.
		
	\begin{figure}[H]
		\centering
		\includegraphics[width=0.6\linewidth]{../../../build/lr2_screen/plot_a2.5_b6.5_no_resave}
		\caption{Точка $O_1$ устойчива, параметр $b$ находится на границе устойчивости}
		\label{fig:boundary}
	\end{figure}
		
	При больших значениях параметров система переходит в режим хаоса (рисунок~\ref{fig:high_b}). 
	
	\begin{figure}[H]
		\centering
		\includegraphics[width=0.6\linewidth]{../../../build/lr2_screen/plot_a10.0_b60.0_no_resave}
		\caption{Все три точки неустойчивы}
		\label{fig:high_b}
	\end{figure}
	
	\begin{figure}[H]
		\centering
		\includegraphics[width=0.6\linewidth]{../../../build/lr2_screen/plot_a10.0_b29.0_no_resave}
		\caption{старт вблизи точки $O_2$}
		\label{fig:strange_O2}
	\end{figure}
	
	На рисунке ~\ref{fig:strange_O2} наблюдается характерная «бабочка» Лоренца.
	
	\subsubsection{Исследование точки $O_2$}
	
	\begin{figure}[H]
		\centering
		\includegraphics[width=0.6\linewidth]{../../../build/lr2_screen/plot_a-1.2_b-0.8_no_resave}
		\caption{Неустойчивая точка $O_1$. Точки $O_2$ и $O_3$ устойчивы}
		\label{fig:plota-1}
	\end{figure}
	
	\begin{figure}[H]
		\centering
		\includegraphics[width=0.6\linewidth]{../../../build/lr2_screen/plot_a2.5_b6.5}
		\caption{Все точки неустойчивы}
		\label{fig:O2_boundary}
	\end{figure}
	
	\subsubsection{Исследование точки $O_3$}
	
	\begin{figure}[H]
		\centering
		\includegraphics[width=0.6\linewidth]{../../../build/lr2_screen/plot_a2.5_b8.0_O3_no_resave}
		\caption{Все точки неустойчивы}
		\label{fig:plota2}
	\end{figure}
	
	\begin{figure}[H]
		\centering
		\includegraphics[width=0.6\linewidth]{../../../build/lr2_screen/plot_a10.0_b29.0O3}
		\caption{старт вблизи точки $O_3$, все точки неустойчивы}
		\label{fig:strange_O3}
	\end{figure}
	
	На рисунках~\ref{fig:strange_O2} и~\ref{fig:strange_O3} показан классический странный аттрактор Лоренца при параметрах $a=10.0$, $b=29.0$. Траектории образуют характерную структуру «бабочки».
	
	\newpage
	\section{Вывод}
	
	В ходе выполнения лабораторной работы получен навык численного исследования динамики нелинейной диссипативной динамической системы, обладающей странным аттрактором. 
	
	Определены области изменения параметров a и b, в которых данная динамическая система является диссипативной. Проведен поиск стационарных точек. 
	
	Найдены параметры a и b, при которых в системе существует странный аттрактор, а также получены различные виды динамики системы. При классических параметрах $a=10.0$, $b=28.0$ зафиксирован режим с характерной структурой «бабочки» Лоренца. 
	
	Для расчета траекторий системы при различных параметрах a и b на языке программирования С++ реализован метод Рунге-Кутта четвертого порядка.
	
	\newpage
	\section{Листинг}
	
	\begin{lstlisting}[caption=main.cpp]
		#include <cmath>
		#include <fstream>
		#include <iostream>
		#include <vector>
		
		using namespace std;
		
		class Runge {
			private:
			double a, b;
			double step;
			double x0, y0, z0;
			int num_steps;
			vector<double> t, x, y, z;
			
			double fx(double x, double y, double z) const { return a * (y - x); }
			
			double fy(double x, double y, double z) const { return b * x + y - x * z; }
			
			double fz(double x, double y, double z) const { return -8.0 / 3.0 * z + x * y; }
			
			public:
			Runge(double a, double b, double step, double x0, double y0, double z0, int steps)
			: a(a), b(b), step(step), x0(x0), y0(y0), z0(z0), num_steps(steps) {}
			
			void run() {
				t.clear();
				x.clear();
				y.clear();
				z.clear();
				
				double cur_t = 0.0;
				double cur_x = x0;
				double cur_y = y0;
				double cur_z = z0;
				
				t.push_back(cur_t);
				x.push_back(cur_x);
				y.push_back(cur_y);
				z.push_back(cur_z);
				
				for (int i = 0; i < num_steps; ++i) {
					double k_1_x = fx(cur_x, cur_y, cur_z);
					double k_1_y = fy(cur_x, cur_y, cur_z);
					double k_1_z = fz(cur_x, cur_y, cur_z);
					
					double k_2_x =
					fx(cur_x + step / 2 * k_1_x, cur_y + step / 2 * k_1_y, cur_z + step / 2 * k_1_z);
					double k_2_y =
					fy(cur_x + step / 2 * k_1_x, cur_y + step / 2 * k_1_y, cur_z + step / 2 * k_1_z);
					double k_2_z =
					fz(cur_x + step / 2 * k_1_x, cur_y + step / 2 * k_1_y, cur_z + step / 2 * k_1_z);
					
					double k_3_x =
					fx(cur_x + step / 2 * k_2_x, cur_y + step / 2 * k_2_y, cur_z + step / 2 * k_2_z);
					double k_3_y =
					fy(cur_x + step / 2 * k_2_x, cur_y + step / 2 * k_2_y, cur_z + step / 2 * k_2_z);
					double k_3_z =
					fz(cur_x + step / 2 * k_2_x, cur_y + step / 2 * k_2_y, cur_z + step / 2 * k_2_z);
					
					double k_4_x = fx(cur_x + step * k_3_x, cur_y + step * k_3_y, cur_z + step * k_3_z);
					double k_4_y = fy(cur_x + step * k_3_x, cur_y + step * k_3_y, cur_z + step * k_3_z);
					double k_4_z = fz(cur_x + step * k_3_x, cur_y + step * k_3_y, cur_z + step * k_3_z);
					
					cur_t += step;
					cur_x += step / 6 * (k_1_x + 2 * k_2_x + 2 * k_3_x + k_4_x);
					cur_y += step / 6 * (k_1_y + 2 * k_2_y + 2 * k_3_y + k_4_y);
					cur_z += step / 6 * (k_1_z + 2 * k_2_z + 2 * k_3_z + k_4_z);
					
					t.push_back(cur_t);
					x.push_back(cur_x);
					y.push_back(cur_y);
					z.push_back(cur_z);
				}
			}
			
			void save_to_csv(const string& filename) {
				ofstream file(filename);
				
				file << "#a=" << a << ", b=" << b << endl;
				
				file << "x,y,z" << endl;
				
				for (size_t i = 0; i < t.size(); ++i) {
					file << x[i] << "," << y[i] << "," << z[i] << endl;
				}
				
				file.close();
			}
		};
		
		int main() {
			double step = 0.001;
			int steps = 5000;
			double anti_stationarity_step = 0.1;
			
			{
				double a = -1.2;
				double b = -0.8;
				double val = sqrt(8.0 / 3.0 * (b + 1));
				double x0 = val + anti_stationarity_step;
				double y0 = val + anti_stationarity_step;
				double z0 = b + 1 + anti_stationarity_step;
				cout << "big numbers again" << endl;
				Runge solver(a, b, step, x0, y0, z0, steps);
				solver.run();
				solver.save_to_csv("../labs/lr2/src/strange_attractor.csv");
				system("python ../labs/lr2/src/main.py");
			}
			return 0;
		}
	\end{lstlisting}
	
\end{document}